% Part: first-order-logic
% Chapter: proof-systems
% Section: seqeunt-calculus

\documentclass[../../../include/open-logic-section]{subfiles}

\begin{document}

\iftag{FOL}
      {\olfileid{fol}{prf}{seq}}
      {\olfileid{pl}{prf}{seq}}

\olsection{क्रमवर्ती कलन}

अनेक !!{derivation} पद्धती !!{sentence}s च्या मांडण्यांवर कार्य करतात,
तर क्रमवर्ती कलन \emph{क्रमवर्तींवर} कार्य करते. क्रमवर्ती ही पुढील रूपाची
अभिव्यक्ती असते:
\[
!A_1, \dots, !A_m \Sequent !B_1, \dots, !B_m,
\]
म्हणजे !!{sentence}s च्या अशा दोन अनुक्रमांची जोडी, ज्यांना क्रमवर्ती
चिन्ह~$\Sequent$ वेगळे करते. यांपैकी कोणताही अनुक्रम रिकामा असू शकतो.
क्रमवर्ती कलनातील !!^a{derivation} हा क्रमवर्तींचा वृक्ष असतो. त्यातील
सर्वांत वरच्या क्रमवर्ती विशेष रूपाच्या असतात (त्यांना ``आरंभीच्या
क्रमवर्ती'' किंवा ``स्वयंसिद्धके'' म्हणतात), आणि इतर प्रत्येक क्रमवर्ती
तिच्या लगत वर असलेल्या क्रमवर्तींपासून निगमन नियमांपैकी एखाद्या नियमाने
निष्पन्न होते. निगमन नियम एकतर क्रमवर्तींतील !!{sentence}s हाताळतात
(डाव्या किंवा उजव्या बाजूला त्यांची भर घालतात, त्यांना काढतात किंवा त्यांचा क्रम
बदलतात), किंवा नियमाच्या निष्कर्षात एखादे जटिल !!{formula} आणतात.
उदाहरणार्थ, $\LeftR{\land}$ नियमामुळे $!A, \Gamma \Sequent \Delta$ पासून
$!A \land !B, \Gamma \Sequent \Delta$ हा निष्कर्ष काढता येतो; आणि
$\RightR{\lif}$ नियमामुळे $!A, \Gamma \Sequent \Delta, !B$ पासून
$\Gamma \Sequent \Delta, !A \lif !B$ हा निष्कर्ष काढता येतो. हे कोणत्याही
$\Gamma$, $\Delta$, $!A$ आणि~$!B$ साठी लागू आहे. (विशेषतः, $\Gamma$
आणि~$\Delta$ रिकामे असू शकतात.)

क्रमवर्ती कलनावर आधारित $\Proves$ संबंधाची व्याख्या पुढीलप्रमाणे केली
जाते: $\Gamma \Proves !A$ तेव्हा आणि केवळ तेव्हाच, जेव्हा असा एखादा
अनुक्रम~$\Gamma_0$ अस्तित्वात असतो की $\Gamma_0$ मधील प्रत्येक $!A$
हा~$\Gamma$ मध्ये असतो आणि ज्याच्या मुळाशी $\Gamma_0 \Sequent !A$ ही
क्रमवर्ती आहे असे !!{derivation} अस्तित्वात असते. $\Sequent !A$ या क्रमवर्तीचे
!!a{derivation} असेल तर क्रमवर्ती कलनात $!A$ हे प्रमेय असते. उदाहरणार्थ,
$\Proves (!A \land !B) \lif !A$ हे दाखवणारे !!a{derivation} पुढे दिले आहे:
\begin{prooftree}
  \Axiom$!A \fCenter !A$
  \RightLabel{\LeftR{\land}}
  \UnaryInf$!A \land !B \fCenter !A$
  \RightLabel{\RightR{\lif}}
  \UnaryInf$\fCenter (!A \land !B) \lif !A$
\end{prooftree}

प्रत्येक $!A \in \Gamma_0$ हे~$\Gamma$ मध्ये असेल अशा एखाद्या
$\Gamma_0 \Sequent$ चे !!a{derivation} अस्तित्वात असेल (म्हणजे
क्रमवर्तीची उजवी बाजू रिकामी असेल), तर क्रमवर्ती कलनात संच~$\Gamma$
विसंगत असतो. \RightR{\Weakening} नियम वापरून विसंगत संचापासून कोणतेही
!!{sentence} !!{derive}d करता येते.

क्रमवर्ती कलनाची निर्मिती १९३० च्या दशकात गरहार्ड गेंटझेन यांनी केली.
त्याची रचना पद्धतशीर आणि सममित असल्यामुळे !!{derivation}s ची उपपत्ती
विकसित करण्यासाठी ते अतिशय उपयुक्त आकारिक साधन आहे. क्रमवर्ती कलनात
!!{derivation}s शोधणे तुलनेने सोपे असते; पण ही !!{derivation}s अनेकदा
वाचायला कठीण असतात आणि त्यांचा सिद्धतांशी असलेला संबंधही काही वेळा
सहज दिसत नाही. तरी !!{derivation} पद्धतींकडे पाहण्याची ही अतिशय सुबक
रीत ठरली आहे, आणि अनेक तर्कप्रणालींसाठी क्रमवर्ती-कलन पद्धती उपलब्ध आहेत.

\end{document}
